%% Supplementary Information: a separate document, as Nature Communications requires.
%% It lives at the project root, beside main.tex, so every path is root-relative. Overleaf's
%% BibTeX cannot read files in a parent folder (../references.bib), which is why a root file
%% inside supplementary/ compiled here but lost its references in Overleaf.
%% To compile in Overleaf: Menu > Main document > supplement.tex.
\documentclass[12pt]{article}
\usepackage[margin=2.5cm]{geometry}
\usepackage{graphicx,amsmath,amssymb,booktabs}
\usepackage[numbers,sort&compress]{natbib}
\usepackage{wasysym}
\usepackage{array,longtable,pdflscape}
\usepackage[hidelinks,breaklinks]{hyperref}

% macros that sn-nature.bst writes into the bibliography (defined by sn-jnl in main.tex)
\providecommand{\bibcommenthead}{}
\providecommand{\burl}[1]{\url{#1}}

% Feature-support marks (same macros as main.tex)
\newcommand{\mkfull}{\CIRCLE}
% \mkpart deliberately not defined: the comparison is binary by design.
\newcommand{\mknone}{\Circle}

\emergencystretch=3em

% Supplementary numbering: Section S1, S1.1; Table S1; Figure S1
\renewcommand{\thesection}{S\arabic{section}}
\renewcommand{\thetable}{S\arabic{table}}
\renewcommand{\thefigure}{S\arabic{figure}}

\title{Supplementary Information for\\[4pt]
\large epymorph turns spatial epidemic model comparisons into controlled forecasting experiments}
\author{}
\date{}

\begin{document}
\maketitle

\section{Software landscape systematic search}\label{sec:S1}

%% PLACEHOLDER -- content not yet migrated.
%% Current source of truth (old working folder, not yet ported):
%%   drafts/appendix-software-landscape-v02.md
%% Will also cite the reproducible lit-search pipeline in
%% results/lit-search/ once that migration/condensing step is done.

\section{Supplementary figures and tables for the case studies}\label{sec:S2}

Here we provide additional figures and tables to support our case study analyses. 
Specifically, we demonstrate the rolling forecasts and we report the coverage of each model's prediction intervals. 
Note that across all forecast figures, we display a four-week forecast from each origin week, showing the median, the 50\%, and 95\% prediction intervals.

\textbf{Influenza.} 
Figure~\ref{fig:S2-flu} shows the rolling forecasts of the empirical-mixing model. 
For one example state (Minnesota), we show the forecasts of each age class separately (panel a), and then three example state-level aggregate forecasts (panel b).
Forecasts tended to over-predict admissions during the peak of the season.

\textbf{COVID-19.} 
Figure~\ref{fig:S2-covid} shows the rolling forecasts of the commuting model for the most populous counties in each state (panel a) and for the state-level aggregate data (panel b). The rows show the two fitted quantities: weekly hospital census and deaths.

\textbf{Interval coverage.} 
Figure~\ref{fig:S2-coverage} shows the coverage of each model, i.e., how often observations fell within each model's 95\% prediction interval. The influenza models were generally under-dispersed.
For COVID-19, coverage of the county models' state totals fell across forecast horizons, whereas the state-level model had better coverage near 100\%. 

\textbf{Reproducibility and results by location.}
Every model was run at two independent ensemble seeds.
Tables~\ref{tab:flu-reproducibility} and~\ref{tab:covid-reproducibility} give the main comparisons at each seed.
Note that the ranking of the models was the same at both seeds.
Tables~\ref{tab:flu-by-jurisdiction} and~\ref{tab:covid-by-state} give each model's mean log WIS in every jurisdiction.

\begin{figure}[htbp]
\centering
\includegraphics[width=\textwidth]{figures/figS-flu-rolling-forecasts.pdf}
\caption{Rolling four-week forecasts of weekly influenza hospital admissions from the
empirical-mixing model, 2025--2026 season. (a) Each age class in Minnesota. (b) Total
admissions in Minnesota, Alabama and Washington, summed from the age classes within each
realization. Circles, observations; lines, forecast medians; shading, 50\% and 95\% prediction
intervals.}\label{fig:S2-flu}
\end{figure}

\begin{figure}[htbp]
\centering
\includegraphics[width=\textwidth]{figures/figS-covid-rolling-forecasts.pdf}
\caption{Rolling four-week forecasts of the weekly COVID-19 hospital census and weekly deaths
from the commuting model, 2020--2021 winter wave. (a) The most populous county of each state.
(b) State totals, summed from the 32 counties within each realization. Each county in (a) lies
in the state directly below it. Circles, observations; lines, forecast medians; shading, 50\%
and 95\% prediction intervals.}\label{fig:S2-covid}
\end{figure}

\begin{figure}[htbp]
\centering
\includegraphics[width=\textwidth]{figures/figS-calibration-both-demos.pdf}
\caption{Coverage of 95\% prediction intervals by forecast horizon, averaged over forecast
origins and locations. (a) Influenza, total admissions per jurisdiction. (b) COVID-19, hospital
census and weekly deaths at the state scale. Dotted lines, nominal coverage.}\label{fig:S2-coverage}
\end{figure}

\input{tables/flu-reproducibility}
\input{tables/covid-reproducibility}
% The long table is set in place, not floated: clear the pending floats first so the
% tables print in the order they are numbered.
\clearpage
\input{tables/flu-by-jurisdiction}
\input{tables/covid-by-state}


\bibliographystyle{sn-nature}
\bibliography{references}

\end{document}
